% Gedeeld door de syllabi van labo 3, 4 en 5. Pas dit hier aan, niet in de labobestanden.
\section{De dataset}

\noindent
\rowcolors{2}{tint}{white}
\begin{tabularx}{\textwidth}{@{}L{3.1cm}L{3.2cm}Y@{}}
\toprule
\textbf{Bestand} & \textbf{Omvang} & \textbf{Inhoud en gekende problemen} \\
\midrule
\code{athletes.csv}   & 43 rijen, 7 kolommen   & ID, naam, geboortedatum, geslacht, positie, lengte, gewicht. 13 schrijfwijzen voor 4 posities, gewicht als tekst met een decimale komma (\code{"57,5"}), één atleet twee keer ingevoerd met twee ID's, enkele ontbrekende waarden. \\
\code{tests.csv}      & 114 rijen, 8 kolommen  & Sprint over 10\,m en 30\,m, spronghoogte (countermovement jump) en Yo-Yo-niveau, op 3 testmomenten. Sommige sprinttijden staan in milliseconden; enkele ontbrekende waarden. \\
\code{hr\_session.csv} & 12\,000 rijen, 5 kolommen & Hartslag en snelheid per seconde van 5 atleten tijdens één training van 40 minuten. \\
\bottomrule
\end{tabularx}
\rowcolors{1}{}{}
